\section{A smaller arithmetic Tate cutoff with the source ranks unchanged}
\label{sec:cubic-tate-cutoff-transport}

\paragraph{Status and precise scope.}
This is a conditional-interface proof; checks are recorded in
\texttt{CHECK\_EVENTS.tsv}.
It changes the arithmetic cutoff $\lfloor x^{401}\rfloor$ to
$\lfloor x^{399}\rfloor$ for every $0<\varepsilon<1/256$.
The actual source ranks, hosts, geometric
constant $C_t=401$, interpolation order $\lceil x^{377}\rceil$,
ordinary and mixed local conditions, and analytic realizations are
unchanged.  No numerical upper bound for an elliptic-unit index is
proved.  The cubic mixed theta/Selmer identification and its analytic
realization, explicitly retained in AC-K27, remain inputs.  The rational
Dirichlet theorem \texttt{zeroFreeAll\_of\_topG} is not applied to a
cubic character.

All Lean source locators below refer to commit
\texttt{7cbf7a0c32fca1badf05931879a37c3f8a2b03b4}, under \texttt{src/}.
The argument is a mathematical proof about the displayed definitions;
no Lean execution is asserted.

\paragraph{Notation and the arithmetic gate.}
Put $K=\mathbb Q(\omega)$ and use the actual fields $F_u$ and
$H_u=F_u(\sqrt[3]{269})$, group $G_u$, coefficient rings $O_\ell$,
degrees $b_\ell$, and derived indices $i_{\ell,\eta,e}$ of AC-K25--27.
For a real $C\geq 5$ define
\[
 T_C(u)=\lfloor u^C\rfloor,\qquad
 I_C(H_u,u)=\sum_{1\leq e<T_C(u)}\sum_\ell
 \frac{\log\#k_{O_\ell}}{b_\ell}
 \sum_{\eta\in\widehat G_u}\max(0,i_{\ell,\eta,e}).
 \tag{TC.1}
\]
Every coefficient prime and every character of the full growing group
occurs.  These sums are not restricted to the coefficient $269$, to
selected residual characters, or to a sparse set of Tate weights.

Here is the exact window convention used for the replacement.
For every fixed finite datum in the retained cubic application,
every $0<\varepsilon<1/256$, and every page $h(u)$ satisfying
\[
 u^{2+8\varepsilon}\leq 2^{h(u)}
                   \leq2u^{2+8\varepsilon}\quad(u\geq1),
 \tag{TC.2}
\]
there must be $J\geq1$, $X_0$, and a positive relative slack
$\sigma\in(0,1)$, independent of $X,q,d,u,\ell,\eta,e$, such that:
for every $X\geq X_0$ and every prescribed pair
$\operatorname{WinPair}(\varepsilon,X^{1/J},q,d)$, there is one
$W$ with $X^{1/J}\leq W$ and $W^2\leq X$ for which
\[
 I_{399}(H_u,u)\leq(1-\sigma)\exp(u^\varepsilon/2)
                 \quad\hbox{for every }u\in[W,W^2].
 \tag{TC.3}
\]
The fields in this display are the actual fields on that page and
that one pair $q,d$.  If the retained application has several
simultaneously required fields or finite labels, (TC.3) is imposed
on all of them on the same $W$; this proof does not construct a
common window from separate label-dependent windows.  As in the
source, fixed-datum and $\varepsilon$ dependence of $J,X_0,\sigma$
is allowed.  Alternatively one can assume a little-o estimate,
uniform on these same windows, which supplies any chosen fixed slack.

\paragraph{The retained cubic input, stated without a cutoff.}
Denote by $\mathcal R_{\rm cub}$ the unchanged mixed realization
required after AC-K27.  It supplies the actual source-compatible
families, mixed theta data, their strict-Selmer identification with
the local and coefficient normalizations of AC-K25, and their
analytic comparison to the cubic function in AM-H04.  Its arithmetic
interface at a datum $D$ is indexed by the actual integer $t=D.R.t$:
finiteness and the all-prime sum for every $1\leq e<t$, followed by
the source's one-way mixed analogue of \texttt{LinkN\_at}.
In particular the identified aggregate is
\[
 v_D(\ell)=\log\ell\sum_{1\leq e<D.R.t}
                         \operatorname{shaLen}_D(\ell,e),
 \tag{TC.4}
\]
with its prescribed comparison to the actual strict groups of $H$.
The normalization and every finite comparison cost in that
identification are retained, with the positive budget slack they
require.  No reverse theta inequality is assumed.

This formulation does not assert that $\mathcal R_{\rm cub}$ has
been established.  It states explicitly the remaining input needed
for the cubic conclusion.  The upper envelope $x^{401}$ is not a
parameter of this input: it was used to supply (TC.4) for the actual
$t$.  This finite range is the literal range in \texttt{Data.VV},
\texttt{SizeData\_2.VV}, \texttt{ShaFinOdd}, \texttt{vvTerm},
\texttt{LinkV\_at}, and \texttt{HypV}.  The complete proof of
\texttt{vvRecord} uses \texttt{NData.sha\_eq} only for $0<e<D.R.t$;
\texttt{divinN\_at} then consumes the sum (TC.4), with no extension
of its weight range.  Thus an unproved mixed replacement of these
identifications remains unproved at exactly the same places.
If a proposed cubic realization uses additional Tate weights,
it does not satisfy this stated retained interface, and the present
conditional conclusion does not cover it.

\begin{theorem}[AC-K35: conditional cutoff transport, arithmetic cutoff $399$]
\label{AC-K35}
For the actual source rank construction, the numerical premise
(TC.3), with the unchanged common-window and slack conventions,
supplies the same all-prime strict bound at every actual $D.R.t$
as the former cutoff-$401$ premise.  It supplies the same
\texttt{FReq.req}/\texttt{FReq.At} arithmetic property on the same
host through all source point/family transports.

Consequently, retaining $\mathcal R_{\rm cub}$ and all its other
geometric, automorphic, local, coefficient, and analytic hypotheses,
the cutoff-$399$ version of (EB.1) can replace its cutoff-$401$
numerical premise in that cubic conditional implication to AM-H04.4.
This is not an unconditional assertion of AM-H04.4, a proof of
$\mathcal R_{\rm cub}$, or a bound for (TC.1).
\end{theorem}

\begin{proof}
First, the old cutoff-$401$ premise implies (TC.3) with exactly
the same $J,X_0,\sigma,W$, fields, and labels: for $u\geq1$,
$T_{399}(u)\leq T_{401}(u)$ and every index summand is
nonnegative.  No choice of new common windows is needed for
this implication.  The old requirement for every
$0<\varepsilon<1/52$ restricts to every
$0<\varepsilon<1/256$.  The rest of the proof establishes sufficiency
of the smaller premise; it does not assume an estimate for the
discarded weights $T_{399}(u)\leq e<T_{401}(u)$.

\proofstep{1. A sharp bound for the ranks actually constructed.}
Write
\[
 a=2+8\varepsilon,\quad b=\tfrac92+10\varepsilon,\quad
 c=\tfrac12-\tfrac\varepsilon{16},\quad
 \alpha=\tfrac{797}{2}+\tfrac{929}{16}\varepsilon.
 \tag{TC.5}
\]
The complete definitions in \texttt{PairRowIffQdRows.lean} give
\[
 \begin{gathered}
 x^a\leq n\leq2x^a,\quad f=2n,\quad
 x^c\leq m\leq2x^c,\quad x^b\leq k\leq2x^b,\quad
 r=k+m-1,\quad L=x^\varepsilon,\\
 T=\lceil x^{377}\rceil,\\
 n_W=8\left\lceil\frac{10r}{8}
              \left(f+\frac{f^2(k-1)^3L^2}{m}\right)\right\rceil,
 \\ r+t=8fr+4+Tn_W.
 \end{gathered}
 \tag{TC.6}
\]
These are the actual equalities proved by
\texttt{exists\_ranks\_rankVal\_frame}; $t$ is not being chosen
independently.  Since $x>1$, $b>c$, $m\geq2$, and $k\geq5$,
\[
 f\leq4x^a,\quad r\leq4x^b,\quad
 f+\frac{f^2(k-1)^3L^2}{m}
        \leq132x^{17+(769/16)\varepsilon}.
\]
For the last estimate the two terms have bounds
$4x^a$ and $128x^{2a+3b+2\varepsilon-c}$; the latter exponent
is $17+(769/16)\varepsilon$ and exceeds $a$.
The ceiling contributes less than $8$, so
\[
 n_W\leq5288x^{43/2+(929/16)\varepsilon},\quad
 T\leq2x^{377},\quad 8fr+4\leq132x^\alpha.
\]
Dropping $-r$ from (TC.6) proves
\[
                         t\leq10708x^\alpha.
 \tag{TC.7}
\]
For comparison, $r\geq k$, $f\geq2x^a$, $k-1\geq k/2$,
and $m\leq2x^c$ also give $n_W\geq(5/2)
x^{43/2+(929/16)\varepsilon}$.  Since $T\geq x^{377}$ and
$(8f-1)r+4>0$,
\[
                         t\geq\tfrac52x^\alpha.
 \tag{TC.8}
\]
This confirms that the proof is estimating the actual construction.

Put $D_*=3191/8$.  Uniformly for $0<\varepsilon<1/256$,
\[
 \alpha<\frac{1633185}{4096}<D_*,\qquad
 D_*-\frac{1633185}{4096}=\frac{607}{4096}.
\]
Therefore (TC.7) implies
\[
 x\geq10708^{4096/607}\quad\Longrightarrow\quad t\leq x^{D_*}.
 \tag{TC.9}
\]
In particular $D_*<399$.  No geometric constant, rank, or host was
altered to obtain this reserve.  By (TC.8), an integer cutoff at most
$398$ cannot majorize this same actual rank construction for any fixed
positive $\varepsilon$; no optimality for other rank constructions is
asserted.

\proofstep{2. The comparisons AC-K25--27 at an arbitrary cutoff.}
The proofs of AC-K25--27 hold for any finite positive integer
cutoff $T$.  The per-weight Fourier error in (FC.5) is
$O(D^3\log D_H+D^4\log(2D))$, independent of the weight.
The local and degree-one torsion bounds (FL.4) and (LT.2),
summed over $e<T$, are bounded by a fixed constant times
\[
 D^2T^2\log(8TD)+D\mathcal R T^2+DrT,
 \tag{TC.10}
\]
where $r\leq12$ and
$\mathcal R=O_\varepsilon(x^{\varepsilon/5}(\log x)^2)$.
The host estimates are still
\[
 D=O_\varepsilon(x^{2+8\varepsilon}),\qquad
 \log D_H=O_\varepsilon(x^{2+41\varepsilon/5}(\log x)^2).
\]
For $T=T_C(x)$ with fixed $C\geq5$, these bounds and the exact
determinant specialization (EU.1)--(EU.4) give
\[
 \left|\sum_{1\leq e<T_C(x)}\sum_\ell\log\#S_\ell(H_x,e)
                       -I_C(H_x,x)\right|
 \ll_\varepsilon x^{2C+4+16\varepsilon}(\log x)^2.
 \tag{TC.11}
\]
Indeed the Fourier exponent is $C+8+161\varepsilon/5$,
which is smaller than $2C+4+16\varepsilon$ in this range.
The two other terms of (TC.10) are smaller as well.
Finiteness and nonnegativity are exactly those already proved
for every finite cutoff in AC-K25--27; no infinite signed sum
is introduced.  At $C=399$ the error is
$O_\varepsilon(x^{802+16\varepsilon}(\log x)^2)$.

The implied constants are uniform over the allowed pair $q,d$
and page at that $x$: the displayed discriminant and ramification
bounds use only the defining window inequalities.  A fixed
finite host factor merely affects the fixed implied constant.
For each fixed $\varepsilon>0$ the error is
$o(\exp(x^\varepsilon/2))$.  Thus (TC.3), after increasing
$X_0$, supplies
\[
 \sum_{1\leq e<T_{399}(u)}\sum_\ell\log\#S_\ell(H_u,e)
                  \leq\exp(u^\varepsilon/2)
                         \quad(u\in[W,W^2]).
 \tag{TC.12}
\]
For example the error can be made at most
$(\sigma/2)\exp(u^\varepsilon/2)$ throughout the window.
The normalization by $b_\ell$ and all bad coefficient-prime
Fourier costs are included in (TC.11).

\proofstep{3. The full common-window transport.}
Fix the source finite datum and its constant $c_0\geq1$ in
the relation $n=c_0 2^h$.  Set
\[
                  \xi=c_0^{1/(2+8\varepsilon)}\geq1.
 \tag{TC.13}
\]
This is the same fixed multiplicative rescaling as in the
complete proof of \texttt{WinPairs.lean}; it is not an
independent choice of a field for each weight or coefficient.
Apply the input window at $u$ and put $x=\xi u$, $Z=\xi W$.
Increase its initial threshold so that
\[
 W\geq\xi^3,\qquad W\geq\xi^{3191},\qquad
 \xi W\geq10708^{4096/607},
 \tag{TC.14}
\]
and so that all the original host and pair thresholds hold.
All these requirements depend only on the fixed datum and
$\varepsilon$ and are imposed before choosing the later scale.
For every $x\in[Z,Z^{3/2}]$, one has
\[
 u=x/\xi\geq W,\qquad
 x\leq(\xi W)^{3/2}\leq W^2,\qquad u\leq W^2.
 \tag{TC.15}
\]
The one pair $q,d$ and the page $h(u)$ are retained.
Multiplying (TC.2) by $c_0=\xi^{2+8\varepsilon}$ gives
$x^{2+8\varepsilon}\leq n\leq2x^{2+8\varepsilon}$.
The source pair construction and its monotonicity give all
the unchanged other fields of \texttt{WinDatum}, with common
window exponent $3/2$; the proof uses only (TC.15), not a Tate
cutoff.

For any of the actual \texttt{RankVal} data at this $x$,
(TC.9) and (TC.14) give
\[
 t\leq x^{3191/8}=\xi^{3191/8}u^{3191/8}
                  \leq u^{399}.
 \tag{TC.16}
\]
Since $t$ is integral, $t\leq\lfloor u^{399}\rfloor$.
All summands in (TC.12) are nonnegative, so it bounds the
entire interval $1\leq e<t$, at every coefficient prime.
Also
\[
                 \exp(u^\varepsilon/2)
                       \leq\exp(x^\varepsilon/2)
                       =\sqrt{\exp(x^\varepsilon)}.
 \tag{TC.17}
\]
The finiteness clause for every weight in that interval is
preserved at the same time.  This is precisely the monotonicity
used by \texttt{VVle.mono}; all its finite-prime-set inequalities
follow from the nonnegative full-prime sum.  Algebraic closures
and local coefficient realizations are not selected to depend
on $e$: their canonical strict groups and normalized lengths
are those of the retained identification.  No coefficient prime,
character, or weight is dropped.

One can formalize this replay by replacing the auxiliary
\texttt{WinDatum.bound} field with the resulting bound for
every compatible actual \texttt{RankVal} datum.  All other
fields are identical.  In the complete proof of
\texttt{supplyOfAdmOf.lean}, the bound field is used only in
the last paragraph, where the old $x^{400}$ majorant is
restricted to $R.t$.  Equation (TC.16) performs that last
restriction directly.  The field construction, pair choice,
admissible ranks, \texttt{HostGood}, and all their dependencies
are therefore unchanged.  For the rational source this proves
the pointwise conclusion of \texttt{HostExW} for each
$0<\varepsilon<1/256$ from the corresponding
\texttt{VWinAt(psi,399)} window data.  Literal \texttt{HostExW}
quantifies over every $0<\varepsilon<1/52$; neither that
full-range wrapper nor its \texttt{TopG} application is asserted
here.  For the cubic application the pointwise argument supplies
the strict bound (TC.4) on its actual retained fields, using
(TC.11)--(TC.12), not by pretending those fields are rational
\texttt{IsChiField} instances.

\proofstep{4. Every later interface preserves the actual range.}
In \texttt{Supply.lean}, \texttt{FReq.req} passes $R.t$ to
the arbitrary arithmetic predicate; \texttt{FReq.At} passes
$D.R.t$.  The complete \texttt{FReq.at\_of\_from} proof
uses the actual rank equality, host equality, and the equality
of the selected places to transport that predicate.  It does
not enlarge $t$ or call an arithmetic assertion at exponent
$401$.  \texttt{FamCore.R\_eq} keeps these ranks identical
over the family at the same scale.  Likewise
\texttt{hypV\_of\_cls} is a restriction of the same predicate.
The compositum with the fixed cubic field is included in the
retained structured host data and transported with its fixed
embeddings.  This step is not an arbitrary field isomorphism
which could permute the specified cubic character.

The complete \texttt{vvRecord}/\texttt{divinN\_of\_hypV}
proofs turn the finite-prime-set bounds at $e<D.R.t$ into
the summable quantity (TC.4), then use \texttt{LinkN\_at}.
The latter is an upper bound for the theta quantity in terms
of this sum, with its other geometric terms unchanged.
The argument neither reverses this inequality nor asks for
any $e\geq D.R.t$.  The mixed replacements supplied by
$\mathcal R_{\rm cub}$ enter at exactly these same finite
indices.  Consequently the new arithmetic bound supplies
the same retained family-level conclusion.  All geometric
uses of the fixed constant $C_t=401$, such as the allowed
size of interpolation parameters, remain untouched.

The packaged \texttt{Supply.ExD} has a full-range
\texttt{HostExW} premise, so it cannot simply be invoked with
the restricted supply.  The actual construction has the following
pointwise replay.  In the complete proof of \texttt{hostQ} in
\texttt{SmallestPrimePolynomialBound.lean}, the host supply is
called once, at the fixed $\varepsilon$.  Its \texttt{relevel}
changes the list of level primes but leaves $R$, the underlying
\texttt{HostAx}, and $TT$ unchanged, so the predicate
$G.\mathrm{req}$ and the integer $R.t$ are identical.
The same fixed-$\varepsilon$ argument therefore supplies the
pointwise \texttt{HostExQW} data.

In the complete proof of \texttt{hostWW}, the only such host call
is again at that fixed $\varepsilon$.  If $\varepsilon_0$ and
$\varepsilon_1$ are its existing positive object and weight
thresholds, replace its output threshold by
$\min\{\varepsilon_0,\varepsilon_1,1/512\}$.  Then every
$0<\varepsilon$ below or equal to this threshold is strictly
below $1/256$.  The object and weight constructions, their maxima
of constants, the intermediate kits, and the resulting family
use only this same $\varepsilon$.  Thus the unchanged literal
\texttt{DC.HostWW} conclusion holds, with its smaller existential
threshold; no full-range \texttt{HostExW} is manufactured.

The complete \texttt{kitS} proof then takes the minimum of the
host and object thresholds, and \texttt{designS} takes another
minimum with $1/53$.  Both call their input only at the chosen
$\varepsilon$; \texttt{designS} transports $G.\mathrm{pt}$ through
\texttt{exists\_famP\_of\_kits} into $G.\mathrm{cls}$.
Accordingly the construction used in \texttt{ObjCTrue.recS}
gives the same literal \texttt{DC.DesignChoices} conclusion
with a smaller positive threshold, without applying its packaged
\texttt{ExD} statement.  The complete
\texttt{designExists\_of\_choices} and \texttt{topDesWin} proofs
again only take minima of positive thresholds and evaluate their
inputs at the chosen exponent.  All their other source hypotheses
remain exactly the retained hypotheses.  This proves the family
replay on the smaller interval rather than adding a stronger
family-existence clause to $\mathcal R_{\rm cub}$.

\proofstep{5. The analytic exponent and the cubic scope.}
Steps 1--4 require every $0<\varepsilon<1/256$, which lies
inside the original source interval $0<\varepsilon<1/52$.
This restriction preserves arbitrarily small positive exponents.
Suppose the retained analytic realization gives
its usual window bound $G_{269,j}(y)\leq C y^{1/2+A\varepsilon}$
for arbitrarily large common windows of fixed relative exponent,
with $A\geq0$ fixed as $\varepsilon$ tends to zero, as required by
that realization.  For a possible zero with real part
$\beta>1/2$, choose
\[
 0<\varepsilon<\min\left\{\varepsilon_{\rm realization},
       \frac1{256},\frac1{4\max(1,A)},
       \frac{\beta-1/2}{2\max(1,A)}\right\}.
\]
Then $0<1/2+A\varepsilon<\min(1,\beta)$, and AM-H04's cubic
window criterion applies exactly as before.  If the realization
uses its separately stated exponent function instead, its
unchanged requirement that this function tend to $1/2$
is used in the same way.  No one fixed positive $\varepsilon$
is asserted to exclude every possible zero.

The source \texttt{FamilyClassConnected} also explicitly
quantifies over every sufficiently small positive $\varepsilon$,
and the displayed \texttt{GrowthBoundGivesRiemannHypothesis}
interface has the same convention.  These rational interfaces
verify the quantifier bookkeeping; their conclusion about
Dirichlet characters is not substituted for the unproved cubic
realization.  Since the cubic realization and its inputs were
retained, this proves the stated conditional replacement.
\end{proof}

\paragraph{What the improvement does not accomplish.}
The new premise is weaker than the cutoff-$401$ numerical
premise because it sums a proper shorter initial interval of
nonnegative indices on the same field.  It is still unproved.
Equations (TC.7)--(TC.8) also show that a logarithmic cutoff
cannot dominate the actual source $t$ by this route.
The exact definitions impose $T=\lceil x^{377}\rceil$;
\texttt{ParamsT.TO\_ge} furthermore contains the constraint
\[
 \frac{PP^4 D_p^3 k f x^{C_p}}{m\,\mathrm{de2}}\leq T,
 \qquad PP=C_P(frL)^{e_P},\quad e_P\geq12.
\]
Changing that construction requires a new interpolation estimate.

The higher-K upper envelopes audited in AC-K24 remain too large
at cutoff $399$.  For example, in its 2015 envelope for
$\log\#K$, the factor $D_H^{528m^5d_H^4}$ alone has logarithm
at least a positive constant times $m^5$ when $D_H>1$.
At a weight comparable to $x^{399}$ this is of order
$x^{1995}$, much larger than $x^\varepsilon/2$ for the
same $x$ and $0<\varepsilon<1/256$.  Thus that upper envelope
does not prove the needed $\exp(x^\varepsilon/2)$ estimate.
This is a failure of an available upper bound, not a lower bound
for the actual K-group or elliptic-unit index.

\paragraph{Source and geometric scope.}
The exact new source reads are recorded with hashes in the
adjacent manifest.  The complete rank-construction and rank-bound
proofs, \texttt{WinPairs}, \texttt{supplyOfAdmOf}, \texttt{Supply},
\texttt{DivinNOfHypV}, the complete pointwise chain
\texttt{hostQ}/\texttt{hostWW}/\texttt{kitS}/\texttt{designS},
\texttt{ObjCTrue.recS}, \texttt{designExists\_of\_choices},
\texttt{topDesWin}, and the displayed final analytic interface
were read.  AC-K25--27's applied finite-cutoff comparisons were
rederived; their upstream duality and elliptic-unit theorem
inputs are retained, not freshly reproved here.
The finite arithmetic cover, Artin--Tate coefficients, and
determinant lines are exactly those already described in those
modules.  No new Grothendieck archive search, motive, or
geometric realization is asserted.
